\originalchapter{因果视界、平坦性与暴胀}
本章对应8.902第11讲末至第12讲的暴胀部分。Friedmann方程描述给定组分后怎样演化，却尚未解释为什么早期密度如此均匀、曲率如此小，以及结构的初始扰动从何而来。暴胀是一类早期加速膨胀模型，必须同时说明加速的条件、持续时间及结束机制；它不是在时间轴上增加一个没有动力学解释的箭头。

\originalsection{复合为什么发生在远低于电离能的温度}
宇宙冷却后，反应$p+e\rightleftarrows\mathrm H+\gamma$逐渐形成中性氢。13.6 eV对应约$1.6\times10^5$ K，但每个重子对应大量光子，高能尾部仍可电离氢。热平衡、非简并粒子近似下，令$x=n_e/n_b=n_p/n_b$且忽略氦，Saha关系为
\[
 \frac{x^2}{1-x}=\frac1{n_b}\left(\frac{m_e\kb T}{2\pi\hbar^2}\right)^{3/2}
 \exp\left(-\frac{\chi_{\rm H}}{\kb T}\right),\qquad\chi_{\rm H}=13.6\ \mathrm{eV}.
\]
这里的氢基态、电子自旋统计因子采用常见简化约定；它们不改变指数导致的温度尺度。光子数密度$n_\gamma=2\zeta(3)(\kb T)^3/(\pi^2\hbar^3c^3)$、$\eta_b=n_b/n_\gamma\ll1$。将$n_b=\eta_b n_\gamma$代入，可看出小$\eta_b$要求温度进一步下降才使$x$显著小于1，因而复合温度约为几千K。

\textbf{复合与光子脱耦不是同一个定义}。中性比例涉及化学反应，脱耦涉及散射率$n_e\sigma_Tc$与膨胀率及沿途光学深度的竞争。Saha平衡在实际复合后期失效，精确最后散射面还需非平衡原子过程；$z_*\sim1100$在此作为课程所用尺度，并不是由上述一个平衡式精确算出的观测红移。宇宙微波背景的均匀性因此提出一个因果问题，而非仅是温度计算问题。

\originalsection{三种视界必须分别定义}
从某个初始时刻$t_i$起，光能传播的最大共动径向距离定义\textbf{粒子视界}
\[
 \chi_p(t)=\int_{t_i}^t\frac{c\dd t'}{a(t')},\qquad D_p=a(t)\chi_p(t).
\]
它取决于模型的过去历史和积分是否收敛。\textbf{事件视界}$\chi_e(t)=\int_t^{t_f}c\dd t'/a(t')$则问现在发出的信号将来最多到哪里，只有该未来积分有限时才存在。\textbf{共动Hubble半径}$R_H=c/(aH)$是一个局部膨胀尺度，不一般等于前两者。

若$a\propto t^p$且$0<p<1$，从$t_i=0$积分有
\[
 D_p=\frac{ct}{1-p},\qquad\frac cH=\frac{ct}{p}.
\]
辐射$p=1/2$时$D_p=c/H$；物质$p=2/3$时$D_p=2c/H$。纯de Sitter模型的$H$为常数，有未来事件视界$c/H$，但其粒子视界还取决于选取的过去边界。用$1/(aH)$讨论暴胀的模出入是方便近似，不应将它直接叫作所有模型的因果视界。

复合前已有辐射与物质两组分。忽略曲率和暗能量，可把粒子视界积分精确做到
\[
 \chi_p(a)=\frac{2c}{H_0\Omega_{m,0}}
 \left[\sqrt{\Omega_{r,0}+\Omega_{m,0}a}-\sqrt{\Omega_{r,0}}\right].
\]
这由$H=H_0a^{-2}\sqrt{\Omega_{r,0}+\Omega_{m,0}a}$直接积分得到。复合时的角尺度约为$\vartheta_p=a_*\chi_p(a_*)/D_A(z_*)=\chi_p(a_*)/D_M(z_*)$。选取课程示例参数后它是度数量级，远小于整个天空。原笔记在复合时把辐射与物质之比估为0.05；对$\Omega_{r,0}\simeq9.2\times10^{-5}$、$\Omega_{m,0}=0.315$、$z_*=1090$，该比为$\Omega_{r,0}(1+z_*)/\Omega_{m,0}\simeq0.32$，纯物质近似已有可见误差。

在从减速热大爆炸外推的模型里，大尺度温度相近需要一个初始条件解释：相距远的最后散射区域没有充分的早期因果联系，为什么还具有相近热状态？天空温度近各向同性并非数学上不可能，它指出了模型本身没有产生这一初始条件的机制。暴胀也仍需相应初始区域、势和再加热条件，不能宣称它无条件解决所有初始条件问题。

\originalsection{为什么加速膨胀能改善平坦性与因果联系}
包括真空能在内的总密度参数满足
\[
 \Omega_{\rm tot}-1=\frac{Kc^2}{a^2H^2},\qquad
 \frac{\dd}{\dd t}\left(\frac c{aH}\right)=-\frac{c\ddot a}{\dot a^2}.
\]
因此$\ddot a>0$使共动Hubble半径减小，也压低曲率相对贡献。对恒定$w>-1/3$的平直主导组分，$|\Omega_{\rm tot}-1|\propto a^{1+3w}$：辐射时期随$a^2$增长，物质时期随$a$增长。越向过去外推，平坦性要求越紧；加速期将这个趋势反转。

若暴胀前一片足够大、已具有因果联系的区域被拉伸，后来的可观测区域可来自其内部。扰动波长$\lambda_{\rm phys}=a\lambda_c$增长，而$H$近恒定，于是固定共动模可以从Hubble半径内被拉到外。暴胀结束后的减速阶段$1/(aH)$重新增长，模再进入。在这一叙述中\emph{粒子视界的累计积分仍增长}；减小的是共动Hubble半径，二者不能互换。

无相互作用的重遗迹数密度随$a^{-3}$，$N$个e折使它降低$\ee^{-3N}$。这可稀释在较早相变形成的磁单极，但再加热不能又大量重造同类遗迹。磁单极的产生来自特定大统一对称性破缺，而不是所有量子场论的必然预言。课程中的单极问题是这类模型与观测稀少性之间的张力。

\originalsection{标量场怎样实现暂时的负压强}
本节取$\hbar=c=1$并定义约化Planck质量$\bar M_P^2=1/(8\pi G)$。对负时间号度规，规范标量场作用量为
\[
 S_\phi=\int\sqrt{-g}\dd^4x\left[-\frac12g^{\mu\nu}\partial_\mu\phi\partial_\nu\phi-V(\phi)\right].
\]
均匀场的能量密度和压强为$\epsilon_\phi=\dot\phi^2/2+V$、$p_\phi=\dot\phi^2/2-V$。代入加速方程，忽略其他组分，有$\ddot a>0$当且仅当$\dot\phi^2<V$。缓慢运动的场有接近$-1$的$w$，但仅仅有势能不保证它持续缓慢运动。

变分场作用量，或把上述$\epsilon,p$代入守恒式，得
\begin{equation}\label{eq:astro-inflaton}
 \ddot\phi+3H\dot\phi+V'(\phi)=0,\qquad
 H^2=\frac{\dot\phi^2/2+V}{3\bar M_P^2},\qquad
 \dot H=-\frac{\dot\phi^2}{2\bar M_P^2}.
\end{equation}
$3H\dot\phi$是均匀膨胀导致的阻尼，不是把能量自动转移到热辐射；实际再加热还需场与其他粒子的耦合。直接令$\Lambda$在早期占主导不能说明结束，因为真正恒定的$\Lambda$不会自行衰变。

\originalsection{慢滚条件与持续时间}
当$\dot\phi^2\ll V$且$|\ddot\phi|\ll|3H\dot\phi|$，有
\[
 H^2\simeq\frac V{3\bar M_P^2},\qquad 3H\dot\phi\simeq-V',\qquad
 \epsilon_V=\frac{\bar M_P^2}{2}\left(\frac{V'}V\right)^2,\quad
 \eta_V=\bar M_P^2\frac{V''}V.
\]
$\epsilon_V\ll1,|\eta_V|\ll1$通常保证慢滚近似自洽，精确加速条件则是$\epsilon_H=-\dot H/H^2<1$。将慢滚方程代回去可得$\dot\phi^2/V\simeq2\epsilon_V/3$；原笔记$1/(24\pi G)(V'/V)^2$实际上是这个动能比的近似，和本卷标准$\epsilon_V=1/(16\pi G)(V'/V)^2$不是同一定义。

从起始场值滚到结束场值的e折数为
\begin{equation}\label{eq:astro-efolds}
 N=\ln\frac{a_f}{a_i}=\int_{t_i}^{t_f}H\dd t
 \simeq\frac1{\bar M_P^2}\int_{\phi_f}^{\phi_i}\frac V{V'}\dd\phi.
\end{equation}
减小共动Hubble半径的比值为$(a_iH_i)/(a_fH_f)$，近恒定$H$时约为$\ee^{-N}$。常说的50--60个e折是把某个今日观测尺度接到暴胀和再加热历史后的数量级，并非所有模型统一的固定数。再加热温度、暴胀能标和后续状态方程都会改变它。

\begin{example}
\textbf{二次势给出的慢滚练习。}
取$V=m^2\phi^2/2$，则$\epsilon_V=\eta_V=2\bar M_P^2/\phi^2$。用$\epsilon_V\simeq1$估结束场值$\phi_f\simeq\sqrt2\bar M_P$，积分式\eqref{eq:astro-efolds}得$N\simeq(\phi_i^2-\phi_f^2)/(4\bar M_P^2)$。这是一项解析练习，说明持续时间取决于场值区间，而不能只用势能高度判断。本卷不据它宣称该势符合当前观测。
\end{example}

\originalsection{结构种子、结束与热史的连接}
近de Sitter阶段一个轻的规范标量场的真空涨落在Hubble尺度冻结，典型$\delta\phi\sim H/(2\pi)$。不同区域的场值差意味着滚到同一结束面所需时间差$\delta t\sim-\delta\phi/\dot\phi$，进而产生曲率扰动$\mathcal R\sim-H\delta\phi/\dot\phi$。在单场慢滚、绝热真空与领先阶近似下，其无量纲功率为
\[
 \mathcal P_{\mathcal R}\simeq\frac{H^2}{8\pi^2\bar M_P^2\epsilon_H}.
\]
这里说明物理联系并给出领先阶量子结果；其完整规范扰动量子化超出原第12讲的经典推导，不将这一估计冒充已证明的一般初始谱。多场、非规范动能和非平凡初态均可改变结果。

场接近势的最低点后可能振荡并将能量转移给普通粒子，形成热辐射；暴胀结束、能量转移和达到热平衡可以是不同阶段。以后随膨胀冷却，经历辐射—物质相等、复合、星系形成与晚期暗能量优势。原讲义时间轴的$10^{-34}$秒是示意的模型尺度，而不是测得的统一暴胀时刻。接下来先研究扰动如何在这一背景中增长，再回到可直接由热平衡与反应率计算的早期热史。
